\subsection{Subspaces}

\begin{exercise}{1c1}
    For each of the following subsets of $\fbb^3$, determine whether it is a subspace of $\fbb^3$.
    \begin{subexercises}
        \item $\set{(x_1, x_2, x_3) \in \fbb^3 \mid x_1 + 2x_2 + 3x_3 = 0}$
        \item $\set{(x_1, x_2, x_3) \in \fbb^3 \mid x_1 + 2x_2 + 3x_3 = 4}$
        \item $\set{(x_1, x_2, x_3) \in \fbb^3 \mid x_1x_2x_3 = 0}$
        \item $\set{(x_1, x_2, x_3) \in \fbb^3 \mid x_1 = 5x_3}$
    \end{subexercises}
\end{exercise}

\begin{sol}
    Let $U$ be the subset of $\fbb^3$ in question.
    \begin{subexercises}
        \item Since $0 + 2(0) + 3(0) = 0$, $(0, 0, 0) \in U$, so $U$ is nonempty. Let $u = (u_1, u_2, u_3), v = (v_1, v_2, v_3) \in U$ and $\alpha \in \fbb$.
        Then
        \[
            (u_1 + v_1) + 2(u_2 + v_2) + 3(u_3 + v_3) = (u_1 + 2u_2 + 3u_3) + (v_1 + 2v_2 + 3v_3) = 0 + 0 = 0,
        \]
        and
        \[
            \alpha u_1 + 2(\alpha u_2) + 3(\alpha u_3) = \alpha(u_1 + 2u_2 + 3u_3) = \alpha(0) = 0.
        \]
        Hence, $u + v \in U$ and $\alpha u \in U$, so $U$ is a subspace of $\fbb^3$.
        \item Since $0 + 2(0) + 3(0) = 0 \neq 4$, $(0, 0, 0) \notin U$, so $U$ does not contain the zero vector. Therefore, $U$ is not a subspace of $\fbb^3$.
        \item Since $1 \cdot 0 \cdot 0 = 0$ and $0 \cdot 1 \cdot 1 = 0$, $(1, 0, 0), (0, 1, 1) \in U$.
        However, $(1, 0, 0) + (0, 1, 1) = (1, 1, 1)$ and $1 \cdot 1 \cdot 1 = 1 \neq 0$, so $(1, 1, 1) \notin U$.
        Hence, $U$ is not closed under addition, so $U$ is not a subspace of $\fbb^3$.
        \item Since $5(0) = 0$, $(0, 0, 0) \in U$, so $U$ is nonempty. Let $u = (u_1, u_2, u_3), v = (v_1, v_2, v_3) \in U$ and $\alpha \in \fbb$.
        Then
        \[
            u_1 + v_1 = 5u_3 + 5v_3 = 5(u_3 + v_3),
        \]
        and
        \[
            \alpha u_1 = \alpha(5u_3) = 5(\alpha u_3).
        \]
        Hence, $u + v \in U$ and $\alpha u \in U$, so $U$ is a subspace of $\fbb^3$. \qedhere
    \end{subexercises}
\end{sol}

\begin{exercise}{1c2}
    Verify all assertions about subspaces in Example 1.35.
\end{exercise}

\begin{sol}
    Let $U$ be the subsets in question.
    \begin{subexercises}
        \item \rightimp Suppose $U$ is a subspace of $\fbb^4$.
        Then it must contain the zero vector $(0, 0, 0, 0)$.
        Then $0 = 5(0) + b$, which implies $b = 0$.

        \leftimp Suppose $b = 0$.
        Then $U = \set{(x_1, x_2, x_3, x_4) \in \fbb^4 \mid x_3 = 5x_4}$.
        The proof is similar to the proof in \exref{1c1}(d), so we conclude that $U$ is a subspace of $\fbb^4$.
        \item Since the $0$ function is continuous, $0 \in U$.
        Let $f, g \in U$ and $\alpha \in \rbb$.
        Since both $f + g$ and $\alpha f$ are continuous over $[0, 1]$, $f + g, \alpha f \in U$.
        Therefore, $U$ is a subspace of $\rbb^{[0, 1]}$.
        \item Since the $0$ function is differentiable on $\rbb$, $0 \in U$.
        Let $f, g \in U$ and $\alpha \in \rbb$.
        Since both $f + g$ and $\alpha f$ are differentiable on $\rbb$, $f + g, \alpha f \in U$.
        Hence, $U$ is a subspace of $\rbb^{\rbb}$.
        \item \rightimp Suppose $U$ is a subspace of $\rbb^{(0, 3)}$.
        Then $0 \in U$, so $0'(2) = 0 = b$, which implies $b = 0$.

        \leftimp Suppose $b = 0$.
        Then
        \[
            U = \set{f \in \rbb^{(0, 3)} \mid f'(2) = 0}.
        \]
        The zero function is in $U$ since $0'(2) = 0$.
        Let $f, g \in U$ and $\alpha \in \rbb$.
        Then
        \[
            (f + g)'(2) = f'(2) + g'(2) = 0 + 0 = 0,
        \]
        and
        \[
            (\alpha f)'(2) = \alpha f'(2) = \alpha \cdot 0 = 0.
        \]
        Thus, $f + g, \alpha f \in U$, so $U$ is a subspace of $\rbb^{(0, 3)}$.
        \item The sequence of zeroes has limit 0, so it is in $U$.
        Let $\set{z_n}_{n = 1}^{\infty}, \set{w_n}_{n = 1}^{\infty} \in U$ and $\alpha \in \cbb$.
        Then
        \[
            \lim_{n \to \infty} (z_n + w_n) = \lim_{n \to \infty} z_n + \lim_{n \to \infty} w_n = 0 + 0 = 0,
        \]
        and
        \[
            \lim_{n \to \infty} (\alpha z_n) = \alpha \lim_{n \to \infty} z_n = \alpha \cdot 0 = 0.
        \]
        Thus, $\set{z_n}_1^{\infty} + \set{w_n}_1^{\infty}, \alpha\set{z_n}_1^{\infty} \in U$, so $U$ is a subspace of $\cbb^{\infty}$. \qedhere
    \end{subexercises}
\end{sol}

\begin{exercise}{1c3}
    Show that the set of differentiable real-valued functions $f$ on the interval $(-4, 4)$ such that $f'(-1) = 3f(2)$ is a subspace of $\rbb^{(-4, 4)}$.
\end{exercise}

\begin{sol}
    Let $U = \set{f \in \rbb^{(-4, 4)} \mid f'(-1) = 3f(2)}$.
    Since the zero function is differentiable on $(-4, 4)$ and $0'(-1) = 0 = 3(0) = 3 \cdot 0(2)$, $0 \in U$, so $U$ is nonempty.
    Let $f, g \in U$ and $\alpha \in \rbb$.
    Then
    \[
        (f + g)'(-1) = f'(-1) + g'(-1) = 3f(2) + 3g(2) = 3(f(2) + g(2)) = 3(f + g)(2),
    \]
    and
    \[
        (\alpha f)'(-1) = \alpha f'(-1) = \alpha(3f(2)) = 3(\alpha f(2)) = 3(\alpha f)(2).
    \]
    Hence, $f + g, \alpha f \in U$, so $U$ is a subspace of $\rbb^{(-4, 4)}$. \qedhere
\end{sol}

\begin{exercise}{1c4}
    Suppose $b \in \rbb$.
	Show that the set of continuous real-valued functions $f$ on the interval $[0, 1]$ such that
    \[
    	\int_0^1 f = b
    \]
    is a subspace of $\rbb^{[0, 1]}$ if and only if $b = 0$.
\end{exercise}

\begin{sol}
    Define
    \[
        U = \set*{f \in \rbb^{[0, 1]} \setmid f \text{ is continuous and } \int_0^1 f = b}.
    \]
    \rightimp Suppose $U$ is a subspace of $\rbb^{[0, 1]}$.
    Then $0 \in U$, so $\int_0^1 0 = 0 = b$, which implies $b = 0$.

    \leftimp Suppose $b = 0$.
    Then
    \[
        U = \set*{f \in \rbb^{[0, 1]} \setmid f \text{ is continuous and } \int_0^1 f = 0}.
    \]
    Since the zero function is continuous on $[0, 1]$ and $\int_0^1 0 = 0$, $0 \in U$, so $U$ is nonempty.
    Let $f, g \in U$ and $\alpha \in \rbb$.
    Then
    \[
        \int_0^1 (f + g) = \int_0^1 f + \int_0^1 g = 0 + 0 = 0,
    \]
    and
    \[
        \int_0^1 (\alpha f) = \alpha \int_0^1 f = \alpha(0) = 0.
    \]
    Therefore, $f + g, \alpha f \in U$, so $U$ is a subspace of $\rbb^{[0, 1]}$.
\end{sol}

\begin{exercise}{1c5}
    Is $\rbb^2$ a subspace of the complex vector space $\cbb^2$?
\end{exercise}

\begin{sol}
    No, $\rbb^2$ is not a subspace of $\cbb^2$.
	Let $x = (1, 0) \in \rbb^2$ and $\alpha = i \in \cbb$.
    Then $\alpha x = (i, 0) \notin \rbb^2$, so $\rbb^2$ is not closed under scalar multiplication and is not a subspace of $\cbb^2$.
\end{sol}

\begin{exercise}{1c6}
    \begin{subexercises}
        \item Is $\set{(a, b, c) \in \rbb^3 \mid a^3 = b^3}$ a subspace of $\rbb^3$?
        \item Is $\set{(a, b, c) \in \cbb^3 \mid a^3 = b^3}$ a subspace of $\cbb^3$?
    \end{subexercises}
\end{exercise}

\begin{sol}
    Let $U$ be the subset in question.
    \begin{subexercises}
        \item Yes, $U$ is a subspace of $\rbb^3$.
        To see this, note that $0^3 = 0^3$, so $(0, 0, 0) \in U$, and $U$ is nonempty.
        Let $u = (u_1, u_2, u_3), v = (v_1, v_2, v_3) \in U$ and $\alpha \in \rbb$.
        Then $u_1^3 = u_2^3$ implies
        \[
            u_1^3 - u_2^3 = (u_1 - u_2)(u_1^2 + u_1u_2 + u_2^2) = 0.
        \]
        If $u_1 - u_2 = 0$, then $u_1 = u_2$ and $(u_1 + v_1)^3 = (u_2 + v_2)^3$.
        If $u_1^2 + u_1u_2 + u_2^2 = 0$, we may complete the square to get
        \[
            u_1^2 + u_1u_2 + u_2^2 = \left(u_1 + \frac{u_2}{2}\right)^2 + \frac{3u_2^2}{4} = 0.
        \]
        Since a square of a real number is nonnegative, we must have $u_1 + \frac{u_2}{2} = 0$ and $\frac{3u_2^2}{4} = 0$, which implies $u_1 = u_2 = 0$.
        In either case, $u_1^3 = u_2^3$ implies $u_1 = u_2$.
        A similar argument shows that $v_1^3 = v_2^3$ implies $v_1 = v_2$.
        Hence,
        \[
            (u_1 + v_1)^3 = (u_2 + v_2)^3,
        \]
        and
        \[
            (\alpha u_1)^3 = \alpha^3 u_1^3 = \alpha^3 u_2^3 = (\alpha u_2)^3.
        \]
        Therefore, $u + v \in U$ and $\alpha u \in U$, so $U$ is a subspace of $\rbb^3$.
        \item No, $U$ is not a subspace of $\cbb^3$.
        To see this, suppose $z^3 = 1$ for some $z \in \cbb$.
        Then $z^3 - 1 = (z - 1)(z^2 + z + 1) = 0$.
        Then $z = 1$ or $z^2 + z + 1 = 0$.
        If $z^2 + z + 1 = 0$, we may use the quadratic formula to obtain
        \[
            z = \frac{-1 \pm \sqrt{1 - 4}}{2} = \frac{-1 \pm i\sqrt{3}}{2}.
        \]
        Let
        \[
            \omega_1 = \frac{-1 + i\sqrt{3}}{2} \quad \text{and} \quad \omega_2 = \frac{-1 - i\sqrt{3}}{2}.
        \]
        Then $\omega_1^3 = \omega_2^3 = 1$, and $\omega_1 + \omega_2 = -1$.
        Let $u = (\omega_1, 1, 0), v = (\omega_2, 1, 0) \in U$.
        Then $u + v = (\omega_1 + \omega_2, 1 + 1, 0 + 0) = (-1, 2, 0)$.
        Since $(-1)^3 = -1 \neq 8 = 2^3$, $u + v \notin U$.
        Hence, $U$ is not closed under addition and is not a subspace of $\cbb^3$. \qedhere
    \end{subexercises}
\end{sol}

\begin{exercise}{1c7}
    Prove or give a counterexample: If $U$ is a nonempty subset of $\rbb^2$ such that $U$ is closed under addition and under taking additive inverses (meaning $-u \in U$ whenever $u \in U$), then $U$ is a subspace of $\rbb^2$.
\end{exercise}

\begin{sol}
    Consider the subset $\qbb^2$ of $\rbb^2$.
	Then $\qbb^2$ is nonempty because $(0, 0) \in \qbb^2$, closed under addition since the sum of any two rational numbers is rational, and closed under taking additive inverses since the additive inverse of a rational number is rational.
	However, for any irrational number $r \in \rbb$, the vector $(r, 0) = r(1, 0) \in \rbb^2$ is not in $\qbb^2$.
	Hence, $\qbb^2$ is not closed under scalar multiplication and is not a subspace of $\rbb^2$.
	Therefore, the statement is false.
\end{sol}

\begin{exercise}{1c8}
    Give an example of a nonempty subset $U$ of $\rbb^2$ such that $U$ is closed under scalar multiplication, but $U$ is not a subspace of $\rbb^2$.
\end{exercise}

\begin{sol}
    Consider the subset
    \[
    	U = \{(x, 0) \mid x \in \rbb\} \cup \{(0, y) \mid y \in \rbb\}.
    \]
    Then $U$ is nonempty since $(1, 0) \in U$, and it is closed under scalar multiplication since any scalar multiple of $(1, 0)$ or $(0, 1)$ is still in $U$.
	However, $U$ is not closed under addition since $(1, 0) + (0, 1) = (1, 1) \notin U$.
	Hence, $U$ is not a subspace of $\rbb^2$.
\end{sol}

\begin{exercise}{1c9}
    A function $f : \rbb \to \rbb$ is called \term{periodic} if there exists a positive number $p$ such that $f(x + p) = f(x)$ for all $x \in \rbb$.
	Is the set of periodic functions from $\rbb$ to $\rbb$ a subspace of $\rbb^{\rbb}$.
	Explain.
\end{exercise}

\begin{sol}
    It is not a subspace of $\rbb^{\rbb}$.
	Consider the periodic functions $f(x) = \cos x$ and $g(x) = \cos \sqrt 2x$.
	Then the sum $f + g$ attains a maximum value of 2 at 0.
	If $f + g$ were periodic with period $p$, then $(f + g)(p) = (f + g)(0) = 2$.
	This implies $\cos p = \cos p\sqrt 2 = 1$.
	Then $p = 2\pi m$ and $p\sqrt 2 = 2\pi n$ for positive integers $m$ and $n$.
	Since $p \sqrt 2 = 2 \pi n$ implies $p = \sqrt 2 \pi n$, we have
    \[
    	2\pi m = \sqrt 2\pi n \implies \sqrt 2 = \frac nm,
    \]
    which is a contradiction, since $\sqrt 2$ is irrational.
	Then $f + g$ is not periodic.
\end{sol}

\begin{exercise}{1c10}
    Suppose $V_1$ and $V_2$ are subspaces of $V$.
	Prove that the intersection $V_1 \cap V_2$ is a subspace of $V$.
\end{exercise}

\begin{sol}
    Let $U = V_1 \cap V_2$.
	Since both $V_1$ and $V_2$ are subspaces of $V$, they both contain the zero vector of $V$.
	Hence, $0 \in U$, so $U$ is nonempty.
	Let $u, v \in U$.
	Then $u, v \in V_1$ and $u, v \in V_2$.
	Since both $V_1$ and $V_2$ are subspaces, they are closed under addition.
	Therefore, $u + v \in V_1$ and $u + v \in V_2$, which implies that $u + v \in U$.
	Let $a \in \fbb$.
	Then since both $V_1$ and $V_2$ are closed under scalar multiplication, we have that $au \in V_1$ and $au \in V_2$, which implies that $au \in U$.
	Therefore, $U$ is a subspace of $V$.
\end{sol}

\begin{exercise}{1c11}
    Prove that the intersection of every collection of subspaces of $V$ is a subspace of $V$.
\end{exercise}

\begin{sol}
    Let $\ccal$ be a collection of subspaces of $V$, and let
    \[
    	U = \bigcap_{W \in \ccal} W.
    \]
    Since each subspace in $\ccal$ contains the zero vector, $0 \in U$, so $U$ is nonempty.
	Let $u, v \in U$.
	Then $u, v \in W$ for all $W \in \ccal$.
	Since each $W$ is a subspace, they are closed under addition, so $u + v \in W$ for all $W \in \ccal$, which implies $u + v \in U$.
	Let $a \in \fbb$.
	Then $au \in W$ for all $W \in \ccal$, which implies $au \in U$.
	Therefore, $U$ is a subspace of $V$.
\end{sol}

\begin{exercise}{1c12}
    Prove that the union of two subspaces of $V$ is a subspace of $V$ if and only if one of the subspaces is contained in the other.
\end{exercise}

\begin{sol}
    Let $U, W$ be subspaces of $V$.
	Suppose $U \subseteq W$.
	Then $U \cup W = W$, which is a subspace of $V$.
	Similarly, if $W \subseteq U$, then $U \cup W = U$, which is a subspace of $V$.

    Conversely, suppose $U \cup W$ is a subspace of $V$.
	If neither $U \subseteq W$ nor $W \subseteq U$, then there exist $u \in U \setminus W$ and $w \in W \setminus U$.
	Since $U \cup W$ is a subspace, it must be closed under addition.
	Then $u + w \in U \cup W$.
	However, $u + w \notin U$ because if it were, then $w = (u + w) - u \in U$, which is a contradiction.
	Similarly, $u + w \notin W$ because if it were, then $u = (u + w) - w \in W$, which is also a contradiction.
	Hence, $u + w \notin U \cup W$, which contradicts the assumption that $U \cup W$ is a subspace.
	Therefore, one of the subspaces must be contained in the other.
\end{sol}

\begin{exercise}*{1c13}
    Prove that the union of three subspaces of $V$ is a subspace of $V$ if and only if one of the subspaces contains the other two.

    \note{This exercise is surprisingly harder than \exref{1c12}, possibly because this exercise is not true if we replace $\fbb$ with a field containing only two elements.}
\end{exercise}

\begin{soliff}
    \rightimp Let $U_1, U_2, U_3$ be subspaces of $V$.
	Without loss of generality, suppose $U_2 \subseteq U_1$ and $U_3 \subseteq U_1$.
	Then $U_1 \cup U_2 \cup U_3 = U_1$, which is a subspace of $V$.

    \leftimp Conversely, suppose $U_1 \cup U_2 \cup U_3$ is a subspace of $V$.
	Suppose that some pair satisfies containment, say $U_3 \subseteq U_2$.
	Let $W = U_2 \cup U_3 = U_2$.
	Then $U_1 \cup W = U_1 \cup U_2 \cup U_3$.
	By \exref{1c12}, either $U_1 \subseteq U_2$ or $U_2 \subseteq U_1$.
	If the former occurs, then $U_1 \cup U_2 = U_2$, which contains $U_3$.
	If the latter occurs, then $U_1 \cup U_2 = U_1$.
	Since $U_3 \subseteq U_2$, then $U_3 \subseteq U_1$.
	Therefore, one of the subspaces must contain the other two.

    Now suppose no subspace is contained in another.
	Then there exist $u \in U_2 \setminus U_3$ and $v \in U_3 \setminus U_2$.
	Since $U_1 \cup U_2 \cup U_3$ is a subspace, $u + \alpha v \in U_1 \cup U_2 \cup U_3$ for any nonzero $\alpha \in \fbb$.
	If $u + \alpha v \in U_2$, then $\alpha v = (u + \alpha v) - u \in U_2$, which implies $v \in U_2$, a contradiction.
	If $u + \alpha v \in U_3$, then $u = (u + \alpha v) - \alpha v \in U_3$, a contradiction.
	Then $u + \alpha v \in U_1$ for all nonzero $\alpha \in \fbb$.
	Now choose nonzero $\beta, \gamma \in \fbb$ with $\beta \neq \gamma$.
	Then $u + \beta v, u + \gamma v \in U_1$, so
    \[
    	(u + \beta v) - (u + \gamma v) = (\beta - \gamma)v \in U_1.
    \]
    Since $\beta - \gamma \neq 0$, we have $v \in U_1$.
	Similarly, $(u + \beta v) - \beta v = u \in U_1$.
	Since $u \in U_2 \setminus U_3$ and $v \in U_3 \setminus U_2$ were arbitrary, $U_2 \setminus U_3 \subseteq U_1$ and $U_3 \setminus U_2 \subseteq U_1$.

    Now let $w \in U_2 \cap U_3$.
	Since $U_2 \nsubseteq U_3$, there exists $x \in U_2 \setminus U_3 \subseteq U_1$.
	Then $w + x \in U_2$.
	If $w + x \in U_3$, then $x = (w + x) - w \in U_3$, a contradiction.
	Hence, $w + x \in U_2 \setminus U_3 \subseteq U_1$, so $w = (w + x) - x \in U_1$.
	Therefore, $U_2 \cap U_3 \subseteq U_1$, which gives $U_2 = (U_2 \setminus U_3) \cup (U_2 \cap U_3) \subseteq U_1$ and similarly $U_3 \subseteq U_1$, contradicting our assumption.
	Hence, one of the subspaces must contain the other two.
\end{soliff}

\begin{exercise}{1c14}
    Suppose
    \[
    	U = \set{(x, -x, 2x) \in \fbb^3 \mid x \in \fbb} \quad \text{and} \quad W = \set{(x, x, 2x) \in \fbb^3 \mid x \in \fbb}.
    \]
    Describe $U + W$ using symbols, and also give a description of $U + W$ that uses no symbols.
\end{exercise}

\begin{sol}
    Let $u = (x, -x, 2x) \in U$ and $w = (y, y, 2y) \in W$.
	Then
    \[
    	u + w = (x + y, -x + y, 2x + 2y) = (x + y, y - x, 2(x + y)).
    \]
    Observe that $x + y$ can take any value in $\fbb$, and $y - x$ can also take any value in $\fbb$.
	Hence, we can write
    \[
    	U + W = \set{(a, b, 2a) \in \fbb^3 \mid a, b \in \fbb}.
    \]
    In words, $U + W$ is the set of all vectors in $\fbb^3$ whose third component is twice the first component, and whose second component can be any value in $\fbb$.
\end{sol}

\begin{exercise}{1c15}
    Suppose $U$ is a subspace of $V$.
	What is $U + U$?
\end{exercise}

\begin{sol}
    Let $u_1, u_2 \in U$.
	Then $u_1 + u_2 \in U$ since $U$ is closed under addition.
	Hence, $U + U \subseteq U$.
	Conversely, let $u \in U$.
	Then $u = u + 0$, where $0$ is the additive identity in $U$.
	Since $0 \in U$, we have that $u \in U + U$.
	Hence, $U \subseteq U + U$.
	Therefore, $U + U = U$.
\end{sol}

\begin{exercise}{1c16}
    Is the operation of addition on subspaces of $V$ commutative?
	In other words, if $U$ and $W$ are subspaces of $V$, is $U + W = W + U$?
\end{exercise}

\begin{sol}
    Yes, the operation of addition on subspaces of $V$ is commutative.
	Let $u \in U$ and $w \in W$.
	Then $u + w \in U + W$.
	Because addition in $V$ is commutative, we have that $u + w = w + u$, where $w \in W$ and $u \in U$.
	Hence, $w + u \in W + U$.
	Therefore, every element of $U + W$ is also an element of $W + U$, which implies that $U + W \subseteq W + U$.
	By a similar argument, we can show that $W + U \subseteq U + W$.
	Hence, $U + W = W + U$.
\end{sol}

\begin{exercise}{1c17}
    Is the operation of addition on the subspaces of $V$ associative?
	In other words, if $V_1, V_2, V_3$ are subspaces of $V$, is
    \[
    	(V_1 + V_2) + V_3 = V_1 + (V_2 + V_3)?
    \]
\end{exercise}

\begin{sol}
    Yes, the operation of addition on subspaces of $V$ is associative.
	Let $u_1 \in V_1$, $u_2 \in V_2$, and $u_3 \in V_3$.
	Then $(u_1 + u_2) + u_3 \in (V_1 + V_2) + V_3$.
	Because addition in $V$ is associative, we have that $(u_1 + u_2) + u_3 = u_1 + (u_2 + u_3)$, where $u_1 \in V_1$, $u_2 \in V_2$, and $u_3 \in V_3$.
	Hence, $u_1 + (u_2 + u_3) \in V_1 + (V_2 + V_3)$.
	Therefore, every element of $(V_1 + V_2) + V_3$ is also an element of $V_1 + (V_2 + V_3)$, which implies that $(V_1 + V_2) + V_3 \subseteq V_1 + (V_2 + V_3)$.
	By a similar argument, we can show that $V_1 + (V_2 + V_3) \subseteq (V_1 + V_2) + V_3$.
	Hence, $(V_1 + V_2) + V_3 = V_1 + (V_2 + V_3)$.
\end{sol}

\begin{exercise}{1c18}
    Does the operation of addition on the subspaces of $V$ have an additive identity?
	Which subspaces have additive inverses?
\end{exercise}

\begin{sol}
    Yes, the operation of addition on the subspaces of $V$ has an additive identity.
	The additive identity is the zero subspace $\{0\}$, which contains only the zero vector of $V$.
	For any subspace $U$ of $V$, we have that $U + \{0\} = U$.

    If a subspace $U$ of $V$ had an additive inverse $W$, then we would have that $U + W = \{0\}$.
	However, $U \subseteq U + W$ since $u = u + 0$ for any $u \in U$.
	Hence, $U + W = \{0\}$ implies that $U = \{0\}$.
	Therefore, the only subspace of $V$ that has an additive inverse is the zero subspace $\{0\}$.
\end{sol}

\begin{exercise}{1c19}
    Prove or give a counterexample: If $V_1, V_2, U$ are subspaces of $V$ such that
    \[
    	V_1 + U = V_2 + U,
    \]
    then $V_1 = V_2$.
\end{exercise}

\begin{sol}
    The statement is false.
	Consider the vector space $\rbb^2$ and the subspaces
    \[
    	V_1 = \set{(x, 0) \in \rbb^2 \mid x \in \rbb}, \quad V_2 = \set{(0, y) \in \rbb^2 \mid y \in \rbb}, \quad U = \set{(z, z) \in \rbb^2 \mid z \in \rbb}.
    \]
    Then $V_1 + U = \rbb^2$ and $V_2 + U = \rbb^2$, but $V_1 \neq V_2$.
\end{sol}

\begin{exercise}{1c20}
    Suppose
    \[
    	U = \set{(x, x, y, y) \in \fbb^4 \mid x, y \in \fbb}.
    \]
    Find a subspace $W$ of $\fbb^4$ such that $\fbb^4 = U \op W$.
\end{exercise}

\begin{sol}
    Let
    \[
        W = \set{(0, w_2, 0, w_4) \in \fbb^4 \mid w_2, w_4 \in \fbb}.
    \]
    Suppose $(x_1, x_2, x_3, x_4) \in U \cap W$.
    Since $(x_1, x_2, x_3, x_4) \in U$, we have that $x_1 = x_2$ and $x_3 = x_4$.
    Moreover, since $(x_1, x_2, x_3, x_4) \in W$, we have that $x_1 = 0$ and $x_3 = 0$.
    Hence, $x_1 = x_2 = x_3 = x_4 = 0$, so $U \cap W = \{0\}$.
    Let $(y_1, y_2, y_3, y_4) \in \fbb^4$.
    Then we can write
    \[
        (y_1, y_2, y_3, y_4) = (y_1, y_1, y_3, y_3) + (0, y_2 - y_1, 0, y_4 - y_3),
    \]
    where $(y_1, y_1, y_3, y_3) \in U$ and $(0, y_2 - y_1, 0, y_4 - y_3) \in W$.
    Hence, every $x \in \fbb^4$ can be uniquely written as $x = u + w$ for some $u \in U$ and $w \in W$, so $\fbb^4 = U \op W$.
\end{sol}

\begin{exercise}{1c21}
    Suppose
    \[
    	U = \set{(x, y, x + y, x - y, 2x) \in \fbb^5 \mid x, y \in \fbb}.
    \]
    Find a subspace $W$ of $\fbb^5$ such that $\fbb^5 = U \op W$.
\end{exercise}

\begin{sol}
    Let
    \[
        W = \set{(0, 0, w_3, w_4, w_5) \in \fbb^5 \mid w_3, w_4, w_5 \in \fbb}.
    \]
    Suppose $(x_1, x_2, x_3, x_4, x_5) \in U \cap W$.
    Since $(x_1, x_2, x_3, x_4, x_5) \in U$, we have that $x_3 = x_1 + x_2$, $x_4 = x_1 - x_2$, and $x_5 = 2x_1$.
    Moreover, since $(x_1, x_2, x_3, x_4, x_5) \in W$, we have that $x_1 = 0$ and $x_2 = 0$.
    Hence, $x_1 = x_2 = x_3 = x_4 = x_5 = 0$, so $U \cap W = \{0\}$.
    Let $(y_1, y_2, y_3, y_4, y_5) \in \fbb^5$.
    Then we can write
    \[
        (y_1, y_2, y_3, y_4, y_5) = (y_1, y_2, y_1 + y_2, y_1 - y_2, 2y_1) + (0, 0, y_3 - (y_1 + y_2), y_4 - (y_1 - y_2), y_5 - 2y_1),
    \]
    where $(y_1, y_2, y_1 + y_2, y_1 - y_2, 2y_1) \in U$ and $(0, 0, y_3 - (y_1 + y_2), y_4 - (y_1 - y_2), y_5 - 2y_1) \in W$.
    Hence, every $x \in \fbb^5$ can be written as $x = u + w$ for some $u \in U$ and $w \in W$, so $\fbb^5 = U \op W$.
\end{sol}

\begin{exercise}{1c22}
    Suppose
    \[
    	U = \set{(x, y, x + y, x - y, 2x) \in \fbb^5 \mid x, y \in \fbb}.
    \]
    Find three subspaces $W_1, W_2, W_3$ of $\fbb^5$, none of which equals $\set 0$, such that $\fbb^5 = U \op W_1 \op W_2 \op W_3$.
\end{exercise}

\begin{sol}
    Let
    \begin{align*}
        W_1 &= \set{(0, 0, w_3, 0, 0) \in \fbb^5 \mid w_3 \in \fbb}, \\
        W_2 &= \set{(0, 0, 0, w_4, 0) \in \fbb^5 \mid w_4 \in \fbb}, \\
        W_3 &= \set{(0, 0, 0, 0, w_5) \in \fbb^5 \mid w_5 \in \fbb}.
    \end{align*}
    Consider
    \[
        W = \set{(0, 0, z_3, z_4, z_5) \in \fbb^5 \mid z_3, z_4, z_5 \in \fbb}.
    \]
    It is straightforward to show that $W = W_1 \op W_2 \op W_3$.
    By \exref{1c21}, we have that $\fbb^5 = U \op W$.
    By uniqueness of the decomposition of $v \in \fbb^5$ as $v = u + w$ for some $u \in U$ and $w \in W$, and uniqueness of the decomposition of $w \in W$ as $w = w_1 + w_2 + w_3$ for some $w_1 \in W_1$, $w_2 \in W_2$, and $w_3 \in W_3$, we have that $\fbb^5 = U \op W_1 \op W_2 \op W_3$.
\end{sol}

\begin{exercise}{1c23}
    Prove or give a counterexample: If $V_1, V_2, U$ are subspaces of $V$ such that
    \[
    	V = V_1 \op U \quad \text{ and } \quad V = V_2 \op U,
    \]
    then $V_1 = V_2$.

    \hint{When trying to discover whether a conjecture in linear algebra is true or false, it is often useful to start by experimenting in $\fbb^2$.}
\end{exercise}

\begin{sol}
    See \exref{1c19} for the definitions of $V_1, V_2, U$, and let $V = \rbb^2$.
    Suppose $(x_1, y_1) \in V_1 \cap U$.
    Since $(x_1, y_1) \in U$, we have that $x_1 = y_1$.
    Since $(x_1, y_1) \in V_1$, we have that $y_1 = 0$.
    Hence, $x_1 = y_1 = 0$, so $V_1 \cap U = \{0\}$.
    Similarly, if $(x_2, y_2) \in V_2 \cap U$, then $x_2 = y_2 = 0$, so $V_2 \cap U = \{0\}$.
    Let $(x, y) \in V$.
    Then we can write
    \[
        (x, y) = (x - y, 0) + (y, y),
    \]
    where $(x - y, 0) \in V_1$ and $(y, y) \in U$.
    Additionally, we can write
    \[
        (x, y) = (0, y - x) + (x, x),
    \]
    where $(0, y - x) \in V_2$ and $(x, x) \in U$.
    Hence, $V = V_1 \op U$ and $V = V_2 \op U$, but $V_1 \neq V_2$.
\end{sol}

\begin{exercise}{1c24}
    A function $f : \rbb \to \rbb$ is called \term{even} if
    \[
    	f(-x) = f(x)
    \]
    for all $x \in \rbb$.
	A function $f : \rbb \to \rbb$ is called \term{odd} if
    \[
    	f(-x) = -f(x)
    \]
    for all $x \in \rbb$.
	Let $V_e$ denote the set of real-valued even functions on $\rbb$, and let $V_o$ denote the set of real-valued odd functions on $\rbb$.
	Show that $\rbb^{\rbb} = V_e \op V_o$.
\end{exercise}

\begin{sol}
    Suppose $f \in V_e \cap V_o$.
    Then $f(-x) = f(x)$ and $f(-x) = -f(x)$ for all $x \in \rbb$.
    Hence, $f(x) = -f(x)$ for all $x \in \rbb$, which implies that $f(x) = 0$ for all $x \in \rbb$.
    
    Now let $g \in \rbb^{\rbb}$.
    Define the functions $g_e, g_o \in \rbb^{\rbb}$ by
    \[
        g_e(x) = \frac{g(x) + g(-x)}{2} \quad \text{and} \quad g_o(x) = \frac{g(x) - g(-x)}{2}.
    \]
    Then
    \[
        g_e(-x) = \frac{g(-x) + g(x)}{2} = g_e(x) \quad \text{and} \quad g_o(-x) = \frac{g(-x) - g(x)}{2} = -g_o(x),
    \]
    so $g_e \in V_e$ and $g_o \in V_o$.
    Moreover, we have that
    \[
        g(x) = g_e(x) + g_o(x),
    \]
    so $\rbb^{\rbb} = V_e \oplus V_o$.
\end{sol}